% sec04a-rules.tex — Section 4 subsection: supergraph Feynman rules and conventions
% Writer: 作家_规则 (Writer_Rules). Body LaTeX only; shared macros live in main.tex.

\subsection{Supergraph Feynman rules and diagram conventions}%
\label{sec:superspace:rules}

This subsection fixes the Feynman rules, insertion operators, graph weights, and loop-integral kit from which every one-loop channel of Section~\ref{sec:superspace} is assembled. We write the flat superspace spinor derivatives in the plus/minus frame as $\cD_\pm$, $\bcD_{\dot\pm}$, with
\begin{equation}\label{eq:ss:rules-dsq}
\cD^{2}=2\,\cD_{-}\cD_{+},
\qquad
\bcD^{2}=2\,\bcD_{\dot+}\bcD_{\dot-},
\qquad
\theta_{12}:=\theta_{1}-\theta_{2} .
\end{equation}
The Fourier convention is $e^{ipx}$ with all vertex momenta incoming, so $\BoxE\mapsto-p^{2}$; the overall momentum-conservation factor $(2\pi)^{4}\delta^{4}(p+p')$ is suppressed in every propagator. The adjoint Killing form is $\kappa_{AB}=\delta_{AB}$ throughout.

\subsubsection*{Fermi--Feynman propagators}

In the absorbed-coupling normalization of Section~\ref{sec:gaugebv} the quadratic kernels invert to the Fermi--Feynman propagators
\begin{equation}\label{eq:ss:rules-propV}
\big\langle\cV^{A}(p,1)\,\cV^{B}(-p,2)\big\rangle
=-\frac{2\,\h g^{2}\,\delta^{AB}}{p^{2}}\;\delta^{4}(\theta_{12}) ,
\end{equation}
\begin{equation}\label{eq:ss:rules-propM}
\big\langle\Phi_{r}^{A}(p,1)\,\widetilde\Phi_{s}^{B}(-p,2)\big\rangle
=\delta_{rs}\,\frac{\h g^{2}\,\delta^{AB}}{16\,p^{2}}\;
\bcD_{1}^{2}\cD_{1}^{2}\,\delta^{4}(\theta_{12}) ,
\end{equation}
and the reversed matter orientation uses the opposite chirality projector. The canonical-normalization audits underlying this section instead use the rescaled vector prepotential $u:=\cV/(\sqrt2\,g)$, for which the vector line has the unit propagator $-\h\,\delta^{AB}p^{-2}\delta^{4}(\theta_{12})$; the translation between the two frames is the single map
\begin{equation}\label{eq:ss:rules-canonical}
\cV=\sqrt2\,g\,u ,
\end{equation}
applied to every field, letter, and propagator at once, and no anomalous factor of $g$ or $\sqrt2$ is generated by the change of frame.

\begin{remark}[Status of the propagator slice]\label{rem:ss:rules-slice}
The propagators~\eqref{eq:ss:rules-propV}--\eqref{eq:ss:rules-propM} are the recorded rules on a local, residual-free Fermi--Feynman slice: a unique propagator set, a Nielsen--Kallosh branch, and a complete momentum-space rule set are \emph{not} asserted by the underlying audits, whose stated scope is partial. The downstream computations of this paper use this conditional slice and say so; the complete graph census remains open (Remark~\ref{rem:ss:rules-ghost}).
\end{remark}

\subsubsection*{Vertex inventory}

The one-loop channels use four classes of vertices.
\emph{(a) Cubic gauge words.} The undotted and dotted cubic gauge functionals give the exponent vertices
\begin{equation}\label{eq:ss:rules-hessians}
-\frac{ig}{4\h}\,f^{UCE}\,\cW^{E\gamma}\,
\big(\cD_{C\gamma}-\cD_{U\gamma}\big),
\qquad
+\frac{ig}{4\h}\,f^{UCD}\,\bcW^{D}_{\dot\gamma}\,
\big(\bcD_{C}{}^{\dot\gamma}-\bcD_{U}{}^{\dot\gamma}\big) ,
\end{equation}
where the two slot labels record the marked quantum ports; the second variation already contains the two labeled port assignments, so no $2!$ remains. Two bookkeeping frames appear in the channel computations, and they are never mixed within one graph weight. In the \emph{canonical frame} of~\eqref{eq:ss:rules-canonical} the vertices carry $\mp ig/(4\h)$, the internal lines carry $-\h\delta^{AB}p^{-2}$, and the insertion product is $(-1/(8\sqrt2))(-1/(4\sqrt2))=1/64$. In the \emph{unit-insertion frame} used in Sections~\ref{sec:superspace:seed}--\ref{sec:superspace:bc}, the insertions are the unit-normalized field strengths $\nabla_{+}\cW_{+}$: relative to the canonical frame, each of the two quantum legs of every cubic vertex rescales by $\sqrt2$, so the effective ordered vertex factors are $\mp ig/2$; each insertion rescales by $\sqrt2$, so the insertion product is $1/32$; and each internal line carries the rescaled propagator $-\h/(2p^{2})\,\delta^{4}(\theta_{12})$. The three effects compensate exactly,
\begin{equation}\label{eq:ss:rules-framecomp}
\underbrace{(\sqrt2)^{2}(\sqrt2)^{2}}_{\text{two vertices}}\;
\underbrace{(\sqrt2)(\sqrt2)}_{\text{two insertions}}\;
\underbrace{\Big(\frac12\Big)^{3}}_{\text{three lines}}
=1 ,
\end{equation}
so every channel coefficient is frame-independent. In the unit-insertion frame each ordered pair of cubic gauge vertices therefore contributes
\begin{equation}\label{eq:ss:rules-gaugepair}
\left(+\frac{ig}{2}\right)\left(-\frac{ig}{2}\right)=\frac{g^{2}}{4} ,
\end{equation}
as used in equation~\eqref{eq:ss:vertices}.
\emph{(b) Matter--gauge chain.} In the canonical frame ($u$, $\Phi_{c}=g^{-1}\Phi$) the first link of the matter chain is
\begin{equation}\label{eq:ss:rules-mattervertex}
S^{(1)}_{\mathrm{mat}}
=-\sqrt2\,g\int d^{8}z\;
\widetilde\Phi_{r}^{A}\,\cV^{U}\,(T_{U})^{B}{}_{C}\,\Phi_{r}^{C}
\end{equation}
(in the absorbed frame of Section~\ref{sec:gaugebv} the same link reads $-g^{-2}\int d^{8}z\,\widetilde\Phi\,\cV\,\Phi$),
with the ordered adjoint color chain $(T_{A_{1}}\cdots T_{A_{n}})^{B}{}_{C}=i^{n}f^{A_{1}D_{1}B}f^{A_{2}D_{2}D_{1}}\cdots f^{A_{n}CD_{n-1}}$; the marked matter vertex carries port coefficient $+\sqrt2\,g/\h$ per generator, so the two matter vertices of a marked matter word give $2g^{2}/\h^{2}$, the three propagators give $\h^{3}$, and the ordered adjoint color route gives $f^{ACD}f^{BCE}$.
\emph{(c) Superpotential cubics.} With $W=-\frac{\sqrt2}{6g^{2}}\,\varepsilon_{rst}f^{ABC}\,\Phi_{r}^{A}\Phi_{s}^{B}\Phi_{t}^{C}$, the chiral cubic vertex primitive is
\begin{equation}\label{eq:ss:rules-superpot}
\frac{\sqrt2}{g^{2}}\,\varepsilon_{rst}\,f^{ABC} ,
\end{equation}
with the tilded analog on the antichiral side (the Euclidean weight $\eta_{E}=-1$ is included in~\eqref{eq:ss:rules-superpot}).
\emph{(d) Ghost words.} The BV spectrum contains the ghosts $c$, $\widetilde c$ and the antighost $\bar c$; their role in the accepted channels is stated in the following remark.

\begin{remark}[Ghost and gauge-fixing sector: open census]%
\label{rem:ss:rules-ghost}
No accepted ghost, Nielsen--Kallosh, gauge-fixing, counterterm, or evanescent-mixing graph rule enters any of the twenty-nine nonzero channels of this paper: the accepted anomaly-sector audits record no ghost-loop contribution to those channels, and the complete graph census of these sectors is explicitly flagged open in the underlying audits. Ghost lines appear in the legend of Figure~\ref{fig:ss:legend} for completeness only.
\end{remark}

\subsubsection*{External-letter insertion operators}

A marked external letter leg is not a propagator endpoint: it carries its own insertion operator, while every uncontracted pair of fields is Wick-replaced by~\eqref{eq:ss:rules-propV}--\eqref{eq:ss:rules-propM}. The linearized port operator and its marked derivative are
\begin{equation}\label{eq:ss:rules-port}
K_{+}^{c}:=-\frac{1}{4\sqrt2}\,\cD_{+}\bcD^{2}\cD_{+},
\qquad
A_{(1)}=K_{+}^{c}\,u ,
\end{equation}
\begin{equation}\label{eq:ss:rules-portdminus}
\cD_{-}A_{(1)}
=-\frac{1}{8\sqrt2}\,\cD^{2}\bcD^{2}\cD_{+}\,u ,
\end{equation}
where the second line uses $\cD^{2}=2\cD_{-}\cD_{+}$ of~\eqref{eq:ss:rules-dsq}. Translated by~\eqref{eq:ss:rules-canonical} to the paper's absorbed fields and letters, the linearized field-strength insertion and its semi-chiral derivative are
\begin{equation}\label{eq:ss:rules-portV}
\nabla_{+}\cW_{+}\big|_{\mathrm{lin}}
=g\,K_{+}^{c}\,u
=-\frac{1}{8}\,\cD_{+}\bcD^{2}\cD_{+}\,\cV ,
\qquad
\cD_{-}\nabla_{+}\cW_{+}\big|_{\mathrm{lin}}
=-\frac{1}{16}\,\cD^{2}\bcD^{2}\cD_{+}\,\cV ,
\end{equation}
which sit on the external legs of the letter $b=-(i/\sqrt2)\,\nabla_{+}\cW_{+}$. The external bottom projectors of the matter letters are
\begin{equation}\label{eq:ss:rules-bottom}
\Phi_{r}(p,\theta,\bar\theta)
=e^{\,i\theta p\bar\theta}\,\theta^{+}\,\sqrt2\,\beta_{r}(p),
\qquad
\widetilde\Phi^{r}(q,\theta,\bar\theta)
=e^{-\,i\theta q\bar\theta}\,\gamma^{r}(q) ,
\end{equation}
obeying $\bcD_{\dot\alpha}\Phi_{r}=0$ and $\cD_{\alpha}\widetilde\Phi^{r}=0$.

\begin{remark}[Matter and antichiral ports: what the sources do not name]%
\label{rem:ss:rules-matterport}
The sources define \emph{no} separately named matter insertion operator and \emph{no} antichiral port $K_{-}^{c}$: the same $K_{+}^{c}$ and $\cD_{-}K_{+}^{c}$ of~\eqref{eq:ss:rules-port}--\eqref{eq:ss:rules-portdminus} act on the matter lines of a marked matter word, and the matter letters enter through the bottom projectors~\eqref{eq:ss:rules-bottom}. We do not introduce any operator beyond this recorded content.
\end{remark}

\subsubsection*{Graph weights and Wick symmetry}

Left functional differentiation of a homogeneous ordered word is graded: removing slot $j$ of $X_{1}\cdots X_{n}$ costs $(-1)^{|X_{j}|\sum_{k<j}|X_{k}|}$, and a permutation $\pi$ of labeled legs costs the parity sign $(-1)^{\epsilon(\pi)}$. Every connected coordinate-superspace graph then carries the labeled-graph weight
\begin{equation}\label{eq:ss:rules-weight}
w(G)=\tau_{E}^{|V(G)|}\,\big(-\tau_{E}^{-1}\big)^{|E(G)|}
\Big(\prod_{v}\mathcal C_{v}\Big)
\Big(\prod_{e}K_{e}^{-1}\Big)(-1)^{\epsilon(G)},
\qquad
\tau_{E}=-\frac{1}{\h},\quad -\tau_{E}^{-1}=\h ,
\end{equation}
i.e.\ one factor $\tau_{E}$ per vertex and one factor $\h$ per propagator, with $\mathcal C_{v}$ the ordered vertex coefficient and $K_{e}^{-1}$ the propagator kernel. An unlabeled graph is divided once by the automorphism group preserving field type, chirality, edge orientation, derivative slot, and ordered color word; no general per-graph symmetry-factor table beyond this rule is recorded in the sources. For the mixed-chirality exponential the action-order factor is exactly one,
\begin{equation}\label{eq:ss:rules-wick}
\frac{1}{2!}\left(S_{1}S_{2}+S_{2}S_{1}\right)=S_{1}S_{2},
\end{equation}
and for one fixed directed source attachment the port-preserving Wick factor is likewise one, as in equation~\eqref{eq:ss:wick}. The closed superspace loop is saturated by the identity of Section~\ref{sec:gaugebv},
\begin{equation}\label{eq:ss:rules-saturation}
\delta^{4}(\theta_{12})\,\cD^{2}\bcD^{2}\,\delta^{4}(\theta_{12})
=16\,\delta^{4}(\theta_{12}),
\qquad\text{i.e.}\qquad
\Big[\cD^{2}\bcD^{2}\,\delta^{4}(\theta)\Big]_{\theta=0}=16 ,
\end{equation}
equivalently $\cD^{2}\bcD^{2}(\theta^{2}\bar\theta^{2})=16$ by direct left-Berezin differentiation with~\eqref{eq:ss:rules-dsq}. The two marked source factors of~\eqref{eq:ss:rules-port}--\eqref{eq:ss:rules-portdminus} multiply to
\begin{equation}\label{eq:ss:rules-assembly}
\left(-\frac{1}{8\sqrt2}\right)\left(-\frac{1}{4\sqrt2}\right)=\frac{1}{64},
\qquad
\left(\frac{1}{64}\right)(16)(2)(2)=1 ,
\end{equation}
where the two factors of $2$ are the mixed anticommutators $\{\cD,\bcD\}=2\,\mathsf p$ on the closed loop, $\mathsf p$ the loop momentum. This product is a $D$-algebra factor; it is not four additional Wick contractions.

\subsubsection*{The ordered-triangle integral kit}

Route the ordered triangle by
\begin{equation}\label{eq:ss:rules-routing}
r_{0}=k,\qquad r_{1}=k+q,\qquad r_{2}=k+p+q,\qquad Q:=p+q ,
\end{equation}
(the matter sector uses $r_{1}=k-q$, $r_{2}=k-p-q$). The two-simplex parametrization of the three denominators is
\begin{equation}\label{eq:ss:rules-simplex}
\frac{1}{r_{0}^{2}r_{1}^{2}r_{2}^{2}}
=2\int_{\Delta_{2}}\frac{dx\,dy\,dz\;\delta(1-x-y-z)}{(\ell^{2}+\Delta)^{3}},
\qquad
\ell=k+yq+zQ,
\qquad
\Delta=xy\,q^{2}+xz\,Q^{2}+yz\,p^{2} ,
\end{equation}
with simplex volume $2\int_{\Delta_{2}}1=2\cdot\tfrac12=1$. With $\int_{k}:=\mu^{2\epsilon}\int d^{d}k/(2\pi)^{d}$, $d=4-2\epsilon$, define
\begin{equation}\label{eq:ss:rules-I23}
I_{2}(\Delta):=\int_{\ell}\frac{1}{(\ell^{2}+\Delta)^{2}},
\qquad
I_{3}(\Delta):=\int_{\ell}\frac{1}{(\ell^{2}+\Delta)^{3}},
\qquad
I_{2}(\Delta)\Big|_{1/\epsilon}=\frac{1}{16\pi^{2}\epsilon} .
\end{equation}
Hatted rotational reduction and the master-family identity $\Delta I_{3}=\frac{\epsilon}{2}I_{2}$ of equation~\eqref{eq:ss:Jnidentities} give
\begin{equation}\label{eq:ss:rules-tensor}
\int_{\ell}\frac{\ell^{m}\ell^{n}}{(\ell^{2}+\Delta)^{3}}
=\frac{\widehat\delta^{mn}}{d}\left[I_{2}(\Delta)-\Delta I_{3}(\Delta)\right]
=\frac{\widehat\delta^{mn}}{4}\,I_{2}(\Delta),
\qquad
\frac{1}{4-2\epsilon}\left(1-\frac{\epsilon}{2}\right)=\frac14 ,
\end{equation}
the quarter being exact in $\epsilon$. For the numerator vectors $L_{1}=2k+q$, $L_{2}=2k+p+2q$ the pole part of the ordered triangle is therefore
\begin{equation}\label{eq:ss:rules-pole}
\operatorname*{Pole}_{\epsilon=0}
\int_{k}\frac{L_{1}^{m}L_{2}^{n}}{r_{0}^{2}r_{1}^{2}r_{2}^{2}}
=\frac{\widehat\delta^{mn}}{16\pi^{2}\epsilon} .
\end{equation}
The finite anomaly rests on the evanescent master integral
\begin{equation}\label{eq:ss:rules-master}
\lim_{\epsilon\to0}\,
\mu^{2\epsilon}\!\int\!\frac{d^{\,4-2\epsilon}\ell}{(2\pi)^{4-2\epsilon}}\,
\frac{\mul^{2}}{(\ell^{2}+\Delta)^{3}}
=\frac{1}{32\pi^{2}} ,
\end{equation}
proved as equation~\eqref{eq:ss:masterintegral} of Section~\ref{sec:superspace:cutting}; it is independent of $\Delta$.

\subsubsection*{Line-style legend}

\begin{figure}[t]
\centering
\begin{tikzpicture}
\draw[sV] (0,0) -- (2.4,0);
\node[right, font=\small] at (2.7,0)
  {vector superfield $\cV$, propagator \eqref{eq:ss:rules-propV}};
\draw[sPhi] (0,-0.85) -- (2.4,-0.85);
\node[right, font=\small] at (2.7,-0.85)
  {chiral matter $\Phi_{r}$, propagator \eqref{eq:ss:rules-propM}};
\draw[sPhit] (0,-1.7) -- (2.4,-1.7);
\node[right, font=\small] at (2.7,-1.7)
  {antichiral matter $\widetilde\Phi^{r}$ (reversed arrow)};
\draw[sghost] (0,-2.55) -- (2.4,-2.55);
\node[right, font=\small] at (2.7,-2.55)
  {ghost $c$, $\widetilde c$, $\bar c$ (Remark~\ref{rem:ss:rules-ghost})};
\draw[sleg] (0,-3.4) -- (1.2,-3.4)
  node[sinsert, label=above:{\scriptsize $\nabla_{+}\cW_{+}$}]{}
  -- (2.4,-3.4);
\node[right, font=\small] at (2.7,-3.4)
  {letter insertion: $\nabla_{+}\cW_{+}$, $\nabla_{+}\Phi_{r}$, $\widetilde\Phi^{r}$, $\bcW_{\dot\alpha}$};
\draw[sVcut] (0,-4.35) -- (2.4,-4.35);
\node[font=\scriptsize] at (1.2,-4.35) {$\times$};
\node[right, font=\small] at (2.7,-4.35)
  {Schwinger-cut (Euler-collapsed) line};
\end{tikzpicture}
\caption{Line-style legend for the supergraphs of this section. Wavy lines are vector-superfield propagators; solid lines with forward (reversed) arrows are chiral (antichiral) matter propagators; dashed lines are ghosts. A filled dot on an external leg is a letter insertion carrying its operator, while every uncontracted field pair becomes a propagator. A wavy line with an $\times$ mark is a Schwinger-cut (Euler-collapsed) line of the cutting mechanism of Section~\ref{sec:superspace:cutting}.}
\label{fig:ss:legend}
\end{figure}

The rules of this subsection feed every ordered channel: assembled channel by channel, they produce the seven nonzero families~\eqref{eq:ss:ch-bb}--\eqref{eq:ss:ch-dcb}, each with the common prefactor $-\frac{\h g^{2}}{16\sqrt2\,\pi^{2}}\,f^{ACD}f^{BCE}$, and the fifty-two structural zeros of the census. The comparison of these results with the holomorphic-twist alphabet is deferred to Section~\ref{sec:ht} \cite{BK}.
